% =====================================================================
%  Thu phan hoi diem-theo-diem -- TETC-2026-05-0252
%
%  Tai lieu DOC LAP, bien dich rieng voi main_revision.tex.
%  Tren Overleaf: doi Main document sang file nay khi muon xuat thu.
%
%  Nguyen tac viet:
%    1. Trich nguyen van y reviewer truoc, roi moi tra loi. Khong dien giai
%       lai y ho theo huong de tra loi hon.
%    2. Cho nao KHONG lam duoc thi noi thang la khong lam duoc, kem ly do.
%       Tuyet doi khong viet "we have addressed this" cho mot viec chua lam.
%    3. Moi con so trong thu deu doi chieu duoc bang script trong repo.
%    4. Nam khang dinh phai rut duoc ke ngay o dau thu, khong giau xuong duoi.
%
%  Doi chieu so lieu: python runners/audit_prose.py  (gom ca thu nay)
% =====================================================================
\documentclass[11pt,a4paper]{article}

\usepackage[utf8]{inputenc}
\usepackage[T1]{fontenc}
\usepackage[margin=2.4cm]{geometry}
\usepackage{amsmath,amssymb}
\usepackage{booktabs}
\usepackage{longtable}
\usepackage{array}
\usepackage{enumitem}
\usepackage{xcolor}
\usepackage{titlesec}
\usepackage[hidelinks]{hyperref}

\definecolor{quotegray}{gray}{0.38}
\definecolor{rulegray}{gray}{0.78}
\definecolor{tagink}{HTML}{1F3D6B}   % xanh muc in, du dam de doc khi in den trang

% ---------------------------------------------------------------------
%  Trich y reviewer. Mot vach doc mau xam ben trai + chu nghieng xam: AE
%  doc luot phai phan biet duoc ngay dau la loi ho, dau la loi minh.
%  Boc trong begingroup de mau khong ro ri ra doan tra loi ben duoi.
% ---------------------------------------------------------------------
\newenvironment{revcomment}
  {\par\addvspace{6pt}\begingroup
   \color{quotegray}\itshape\small
   \begin{list}{}{%
     \setlength{\leftmargin}{1.8em}\setlength{\rightmargin}{1.0em}%
     \setlength{\topsep}{0pt}\setlength{\partopsep}{0pt}%
     \setlength{\parsep}{3pt}\setlength{\itemsep}{0pt}}%
   \item[]}
  {\end{list}\endgroup\normalcolor\addvspace{6pt}}

% ---------------------------------------------------------------------
%  Dau muc cho tung y reviewer. Ban truoc dat gach ngang bang \par\nobreak
%  ngay sau tieu de, nen khi tieu de xuong hai dong thi gach chay CAT NGANG
%  giua tieu de (thay ro o AE-1 va AE-4). Bay gio: ma so mau xanh, tieu de
%  in dam, xuong dong roi moi ke gach het chieu rong -- va khoang cach
%  truoc "Response." co dinh nen moi muc deu nhau.
% ---------------------------------------------------------------------
\newcommand{\itemhead}[2]{%
  \par\addvspace{15pt}%
  \noindent{\bfseries\color{tagink}#1}\hspace{0.7em}{\bfseries #2}%
  \par\addvspace{3.5pt}%
  \noindent{\color{rulegray}\rule{\linewidth}{0.5pt}}%
  \par\addvspace{5pt}%
  \nobreak}

\newcommand{\resp}{\par\addvspace{2pt}\noindent\textbf{Response.}\ }
\newcommand{\changed}{\par\addvspace{5pt}\noindent
  {\bfseries\color{tagink}Changes.}\ }

% Tieu de muc: GIU so thu tu -- thu dai 15 trang, AE can co so de dan chieu,
% va chinh trong thu co \ref{sec:selfreported} tro toi mot muc. Bo so la
% cai \ref do in ra rong.
\titleformat{\section}
  {\normalfont\large\bfseries\color{tagink}}{\thesection.}{0.5em}{}
  [\vspace{-5pt}{\color{tagink}\rule{\linewidth}{0.8pt}}]
\titlespacing*{\section}{0pt}{22pt}{10pt}

\setlength{\parskip}{3.5pt}
\setlength{\parindent}{0pt}
\linespread{1.03}

\begin{document}

\begin{center}
  {\Large\bfseries Response to Reviewers}\\[4pt]
  {\normalsize Manuscript TETC-2026-05-0252}\\[2pt]
  {\normalsize\itshape NISQ-Aware Quantum Kernel SVM for Network Intrusion
   Detection: A Regime-Specific Benchmark on NSL-KDD}\\[6pt]
  {\normalsize Minh Tuan Pham, Phuc Hao Do, Nguyen Nang Hung Van,
   Quang Anh Nguyen, Quan Tran Anh Vo}
\end{center}

\vspace{6pt}
{\color{rulegray}\hrule height 0.8pt}
\vspace{10pt}

Dear Editor-in-Chief, Associate Editor, and Reviewers,

We thank the reviewers for reviews that were unusually specific and, in several
places, correct about errors we had not found ourselves. Reviewer~4 identified a
false theorem and a selection step that filters nothing; Reviewer~2 identified
two citation defects, one of which was a reference we could not verify exists;
Reviewer~1 identified an asymmetric tuning protocol that favoured our own
method. Each of these was checked and each was confirmed.

Acting on them changed the paper's conclusion. We ask the editors to read the
following before the point-by-point responses, because it is the most important
thing we have to report.

\section*{Five claims from the submitted version are withdrawn}

Re-running every experiment under a corrected protocol (ten runs instead of
five, symmetric hyper-parameter tuning, a random forest and XGBoost added as
baselines, and results reported on the full KDDTest$^{+}$ partition rather than
a 300-sample subset) does not support five claims made in the submitted
version. We withdraw them rather than qualify them.

\begin{center}
\small
\begin{tabular}{@{}p{7.3cm}p{8.3cm}@{}}
\toprule
\textbf{Submitted version} & \textbf{Measured under the revised protocol} \\
\midrule
Advantage in the low-data regime &
  \emph{Reversed.} Classical baselines win at small $N$; the ordering flips
  between $N=2000$ and $N=5000$ \\
``$+6.7$ points on the rare-attack subset, Cohen's $d=+0.68$'' &
  \emph{Not reproducible.} No rare-class metric was computed by the code that
  produced it \\
Theorem~1: $n^{\ast}=4$ maximises $J$ &
  \emph{False}, as Reviewer~4 states. $J$ is maximised at $n=2$ $(0.551)$ \\
QSVM-ZZ $0.854 >$ SVM-RBF $0.838$ &
  With a tree ensemble in the comparison, XGBoost $0.8503 >$ QSVM-ZZ $0.8469 >$
  random forest $0.8446$, intervals overlapping \\
``Quantum advantage is real and measurable'' &
  Of $110$ controlled comparisons: $21$ favour the quantum kernel, $21$ favour a
  classical baseline, $68$ are inconclusive \\
\bottomrule
\end{tabular}
\end{center}

\vspace{4pt}
We also report three defects that no reviewer raised and that we found while
preparing this revision. They are listed in Sec.~\ref{sec:selfreported} of this
letter and stated in the manuscript itself, in the erratum subsection of
Sec.~III. One of them is an error in the single theoretical result the revision
retains.

What replaces the withdrawn claims is not a weaker version of the same story. It
is a different result: an ordering that \emph{reverses} with training-set size,
a dimension-selection rule that transfers to a second dataset unchanged, and a
measured mechanism, Gram-matrix concentration, that predicts where the
method stops working. We believe this is a more useful contribution than the one
we submitted, and it exists because the reviewers refused the original.

\section*{On verification}

Every number in the manuscript and in this letter is reproducible from the
released repository. Four audit scripts are provided, and we invite the
reviewers to run them:

\begin{center}
\small
\begin{tabular}{@{}llp{8.2cm}@{}}
\toprule
Script & Result & What it checks \\
\midrule
\texttt{audit\_c4.py} & 100/100 &
  Recomputes every published statistic from the raw per-run records using an
  independent implementation, not the pipeline's own functions \\
\texttt{audit\_figures.py} & 36/36 &
  Every figure is newer than its inputs, and every plotted number is re-derived
  from the artifacts \\
\texttt{audit\_prose.py} & 134/134 &
  Every number written in the manuscript text \emph{and in this letter} is
  re-derived from the artifacts \\
\texttt{verify\_lemma1.py} & 15/15 &
  The corrected kernel expansion, and the failure of the submitted one \\
\bottomrule
\end{tabular}
\end{center}

\vspace{4pt}
These audits found four genuine errors in our own revision code before
publication, including one that produced $n^{\ast}=5$ instead of $n^{\ast}=4$.

\vspace{8pt}
{\color{rulegray}\hrule height 0.8pt}

% =====================================================================
\section{Associate Editor}

\itemhead{AE-1}{Advantage claims too generic; evidence limited to few methods
and a small numerical improvement}
\begin{revcomment}
``the reviewers criticize too generic advantage claims versus evidence on a
limited set of methods and by a small numerical improvement''
\end{revcomment}
\resp
We accept this without reservation. The claim of a general advantage has been
withdrawn. The paper is reframed from ``where is there an advantage'' to
``where is a quantum kernel worth its cost'', and the answer is reported as a
regime map in which negative and inconclusive verdicts appear on the same
footing as positive ones: $21$ / $21$ / $68$ over $110$ comparisons within
baseline families, $17$ / $15$ / $78$ under a single Benjamini--Hochberg
correction over all $110$ cells, and tallied separately for the natural-prior
cells on the full test set ($11$ / $9$ / $22$) and the rare-enriched cells
($10$ / $12$ / $46$), since Table~II shows the two protocols are not
interchangeable.
\changed
New Sec.~VI (Regime Map), Fig.~9, and a rewritten abstract and introduction.
The introduction states the withdrawn claims on its first page.

\itemhead{AE-2}{Insufficient literature coverage, including non-existent
references}
\resp
Both citation defects reported by Reviewer~2 were checked independently and both
were confirmed; see R2-5 and R2-6. Four recent studies have been added and the
positioning table rewritten.
\changed
New Table~I (positioning against four benchmarking studies), Sec.~II rewritten,
bibliography corrected.

\itemhead{AE-3}{Reviewer~1 requests validation on an additional dataset}
\resp
UNSW-NB15 has been added under the identical protocol. This is what allows us to
separate findings that transfer from findings that do not: the selection rule
and the entanglement ablation transfer; the sample-complexity reversal does not.
\changed
Sec.~V-F, Fig.~7.

\itemhead{AE-4}{Comparison with non-SVM methods as a prerequisite for practical
advantage claims}
\resp
A random forest and XGBoost were added to every experiment, tuned by the same
procedure as every other model. They change the headline: XGBoost attains the
highest mean macro-$F_1$ at the main operating point.

Two points of scope are now stated in the manuscript so that the comparison is
not read as more than it is. First, every classical baseline consumes the same
four-component representation as the circuit (Sec.~III-B); the paired protocol
isolates the kernel, not the pipeline. Second, a reference arm outside the
paired protocol trains XGBoost and the random forest on all $122$ one-hot
features on the same subsets, seeds and test set (Table~IV). Against
full-feature XGBoost the quantum kernel reaches parity at $N\geq5000$, not a
lead, and full-feature XGBoost on the entire training partition scores $0.804$,
above every four-qubit model in the paper. The practical claim is worded
accordingly.
\changed
All results sections; Figs.~3--8; Sec.~III-B; Sec.~V-D, fourth qualification, and Table~IV.

\itemhead{AE-5}{Reviewer~4 identified issues in the theory}
\resp
Confirmed and acted on. Theorem~1 was false and has been removed with the
objective it maximised; the Pareto stage filtered nothing and has been replaced
by a lexicographic rule with no weights. Two further theory defects that no
reviewer raised are reported in Sec.~\ref{sec:selfreported}.
\changed
Sec.~III rewritten, including an erratum subsection; new Appendix~A.

\itemhead{AE-6}{``NISQ-ready'' not supported by experiments or noisy
simulations}
\resp
The pipeline is re-run under a noise model derived from the calibration data of
a real superconducting backend (\texttt{FakeManilaV2}) carrying gate errors,
readout errors and thermal relaxation. We are explicit that this is
\emph{not} hardware validation: the circuits were not executed on a quantum
processor, and the manuscript now says so in those words. The noisy re-run
covers one of the ten runs on the $300$-record test split; the manuscript
states this, quotes the result to two decimals, and lists the ten-run version
as the next step rather than implying a noisy confidence interval it does not
have.
\changed
Sec.~V-B and Sec.~VII-A, which separates what we measure from what we do not.

% =====================================================================
\section{Reviewer 1}

We are grateful for a review that identified the single methodological flaw
most damaging to our own conclusions (R1-3), and asked the question whose answer
became the paper's main result (R1-7).

\itemhead{R1-1}{Missing QMI 2026 benchmark; Table~I needs 2025--2026 literature}
\begin{revcomment}
``an important recent work is missing: `Benchmarking quantum machine learning
methods for intrusion detection on noisy quantum computers' (Quantum Machine
Intelligence, 2026). This work should be discussed explicitly''
\end{revcomment}
\resp
The work (Cirillo, Esposito and Seo, \emph{Quantum Machine Intelligence},
vol.~8, no.~1, art.~27, 2026) is now cited, read in full, and discussed in
Sec.~II-B, with a row for it in Table~I. We are glad the reviewer insisted: it
is the closest study to ours that we are aware of, and reading it changed what
we say.

Two of its findings corroborate ours from a different direction, and we now say
so explicitly. It compares against SVM, XGBoost, random forest and a neural
network at matched dimensionality and concludes that the classical models retain
higher peak accuracy while the quantum models remain competitive in
low-dimensional feature spaces: the same ordering we now report, with XGBoost
top at our operating point. And it sweeps the qubit count from $2$ to $10$ and
plots $F_1$ against two-qubit gate count, locating good performance in a
$20$--$35$ CNOT band with degradation beyond roughly $60$ CNOTs. Our operating
point is $24$ CNOTs, inside that band, and the collapse we report at $112$ CNOTs
is consistent with their threshold.

The two studies diverge on the mechanism, and we think that divergence is
useful. They attribute the degradation to two-qubit gate error accumulating on
noisy backends. We observe the same collapse under exact, noiseless statevector
simulation, where no gate error exists, and trace it to concentration of the
Gram matrix instead. Hardware noise is therefore not a necessary ingredient; the
bound is intrinsic to the kernel.

What that study does not vary, and what, to our knowledge, no study in this
area varies, is the training-set size. It is fixed there at a single $15\%$
class-balanced subsample under a random $80/20$ split rather than the official
partitions. That is the axis on which our contribution lies.
\changed
Sec.~II-B; Table~I now covers Bowles \emph{et al.} (2024), Schnabel and Roth
(2024), the QMI 2026 benchmark, Carducci (2026), and the concurrent study of
Gillani \emph{et al.} (2026).

\itemhead{R1-2}{Reliance on NSL-KDD; UNSW supplementary was not available}
\resp
UNSW-NB15 is now a first-class part of the study rather than supplementary
material, run under the same protocol with the same audits. Sec.~VII-D also
quantifies a property of both corpora that constrains what either can show:
NSL-KDD has $610$ test rows identical to a training row in both features and
label, and UNSW-NB15 is far more heavily duplicated ($47.3\%$ of training rows
are internal repeats; $25.0\%$ of test rows have an exact duplicate in
training).
\changed
Sec.~IV-A, Sec.~V-F, Sec.~VII-D, Fig.~7.

\itemhead{R1-3}{Asymmetric hyper-parameter protocol: QSVM fixed at $C=1.0$,
classical SVMs validation-selected}
\begin{revcomment}
``The hyperparameter selection protocol is asymmetric $\ldots$ Without such
analysis, it is difficult to determine whether the reported performance
differences are influenced by the selected regularization parameter.''
\end{revcomment}
\resp
The reviewer is right, and the defect ran in our favour. The submitted version
described fixing the quantum regulariser as avoiding bias; it does the opposite,
granting one family a search the other is denied. Every model, quantum and
classical, is now tuned by the same procedure on the same disjoint tuning set of
$2{,}000$ records.

We report two arms rather than one: \emph{tuned per $N$}, in which the search is
repeated at each training-set size, and \emph{frozen}, in which hyper-parameters
are held at the values selected once. Reporting both is what lets us show that
the paper's main new result is not an artefact of re-tuning (see R1-7).
\changed
Sec.~IV-B; Table~III reports both arms across all sizes.

\itemhead{R1-4}{Conclusions overstated; cannot claim general quantum advantage}
\resp
Withdrawn. See the summary table at the head of this letter. The abstract now
leads with the question rather than with a winning number, and reports $21$ /
$21$ / $68$.
\changed
Abstract, Sec.~I, Sec.~VI, Sec.~VIII.

\itemhead{R1-5}{Include XGBoost, CatBoost, TabNet, FT-Transformer, or soften
practical claims}
\resp
We added a random forest and XGBoost, and we did \emph{not} add CatBoost, TabNet
or FT-Transformer. We prefer to state that plainly rather than imply full
coverage.

CatBoost would supply a second gradient-boosting representative without changing
the comparison structure. The deep tabular family would require its own
tuning-budget protocol to be compared fairly, which would enlarge the study from
a controlled quantum-kernel benchmark into a survey of tabular architectures.
Accordingly we have softened the practical claims as the reviewer's alternative
allows, and Sec.~VII-C states explicitly that we do not claim coverage of all
strong tabular methods.
\changed
All experiments; Sec.~VII-C.

\itemhead{R1-6}{NISQ evidence covers only finite-shot error}
\resp
Addressed with a backend-derived noise model; see AE-6. Circuits are transpiled
onto the backend's coupling map and basis set, giving depth $59$ and $44$
two-qubit gates at $n=4$, $r=2$. The entanglement conclusion is unchanged under
noise. As stated under AE-6, this check is a single run and is reported as such.
\changed
Sec.~V-B, Sec.~VII-A.

\itemhead{R1-7}{Is there a crossover point beyond which the classical SVM
becomes preferable?}
\begin{revcomment}
``It would therefore be useful to identify whether there exists a crossover
point beyond which the classical SVM becomes preferable, or clarify why such a
point is not observed.''
\end{revcomment}
\resp
There is a crossover, and it runs in the direction opposite to the one the
question anticipates: the classical baselines are preferable at \emph{small}
$N$, and the quantum kernel becomes preferable at large $N$. Extending the sweep
from $N=1000$ to $N=10^{4}$ under the natural class prior:

\begin{itemize}[leftmargin=1.3em,itemsep=1pt,topsep=3pt]
  \item At $N=100$ the ranking is XGBoost $0.7802$, random forest $0.7695$,
    QSVM-Z $0.7613$, SVM-RBF $0.7296$, linear $0.7261$, QSVM-ZZ $0.6989$.
  \item At $N=10^{4}$ it is QSVM-ZZ $0.7855$, QSVM-Z $0.7787$, SVM-RBF $0.7740$,
    random forest $0.7728$, XGBoost $0.7706$.
  \item The paired difference against XGBoost passes through zero between
    $N=2000$ and $N=5000$ ($-0.0129 \rightarrow +0.0100$) and is
    Holm-corrected-significant at $N=5000$ ($p=0.027$, $d_z=1.22$) and $N=10^{4}$
    ($+0.0149$, $p=0.0078$).
\end{itemize}

Three qualifications are stated in the manuscript and we repeat them here rather
than let the result be read as more than it is. First, the reversal survives
freezing the hyper-parameters: all six baseline--arm combinations change sign in
the same interval, so it is not a re-tuning artefact. Second, it is not uniform
across baselines: against SVM-RBF the difference turns positive but no cell
reaches corrected significance, so the quantum kernel never demonstrably beats
an RBF kernel at any size we tested. Third, whether the reversal is a property
of the natural prior is not decidable with this design, and we say so: the
enriched arm (rare attacks at $10\%$, the sampling our own submitted protocol
used) cannot be run past $N=2000$ without runs sharing samples, and the
crossover occurs between $N=2000$ and $N=5000$, so the enriched arm ends
before the phenomenon begins; within the range the two arms share, the tree
ensembles lead under both priors and the sign agrees in $22$ of $30$ cells.
An earlier draft of this reply, and of the manuscript, stated that enrichment
``removes'' the crossover; that overstated what the data show, and both now
say that enrichment is untested against it. Fourth, the reversal is a statement
about the shared four-component embedding: against XGBoost given all $122$
features the kernel reaches parity at $N\geq5000$, not a lead (Table~IV).

What does follow for the submitted version is that a sweep confined to
$N\leq1000$, under either prior, could not have seen the regime in which the
kernel does best; that, rather than the prior, is why it reported the opposite
conclusion.
\changed
Sec.~V-D, Fig.~6, Tables~III and~IV.

\itemhead{R1-8}{Table~VI ($N=1000$) inconsistent with Table~IV}
\resp
Confirmed, and the discrepancy decomposes into two factors that the submitted
version conflated. Holding $N=1000$ fixed and varying only these:

\begin{center}
\small
\begin{tabular}{@{}llcc@{}}
\toprule
Representation & Test set & QSVM-ZZ & XGBoost \\
\midrule
frozen        & 300-sample subset & 0.8469 & 0.8503 \\
frozen        & full KDDTest$^{+}$ & 0.7959 & 0.8043 \\
refit per $N$ & 300-sample subset & 0.8526 & 0.8666 \\
refit per $N$ & full KDDTest$^{+}$ & 0.8072 & 0.8242 \\
\bottomrule
\end{tabular}
\end{center}

The test set accounts for about $-0.051$ macro-$F_1$ and the representation mode
for about $+0.006$. The revision reports both test sets and states which is
which: Table~II (``Protocol by contribution'') lists, for every figure and
table, the training pool (rare-enriched, as in the submitted protocol, or the
natural prior), the test set ($300$-record split or full KDDTest$^+$), $N$ and
the number of runs, and the captions of Figs.~3--5 repeat it. The same model at
the same $N$ therefore appears with two different scores in the paper, and the
reader can see why.
\changed
Sec.~IV-B and Table~II; \texttt{c4\_table\_iv\_vs\_vi.csv} in the repository.

\itemhead{R1-9}{Classical $F_1$ low relative to values commonly reported on
NSL-KDD}
\resp
The gap is the evaluation protocol, not the models, and it is large enough to be
worth stating explicitly. Using the full $122$ features and the entire training
partition:

\begin{center}
\small
\begin{tabular}{@{}lcc@{}}
\toprule
Protocol & Random forest & XGBoost \\
\midrule
Official split, evaluate on KDDTest$^{+}$ & 0.7765 & 0.8041 \\
Random split \emph{within} KDDTrain$^{+}$ & 0.9978 & 0.9993 \\
\bottomrule
\end{tabular}
\end{center}

Studies reporting $\approx0.99$ on NSL-KDD are, in the cases we could check,
splitting randomly within the training partition rather than using the official
train/test partitions, which places near-duplicate records on both sides. We use
the official split throughout. The four numbers above are now in Sec.~IV-A of
the manuscript, with $0.804$ named as the reference against which every absolute
score should be read.
\changed
Sec.~IV-A; \texttt{c4\_protocol\_vs\_literature.csv} in the repository.

\itemhead{R1-10}{Supplementary material was not accessible}
\resp
The supplementary material is replaced by a public repository containing the
complete implementation, the frozen protocol files, every intermediate artifact,
all figure-generation scripts, and the four audit scripts listed above. The link
and the commit hash are given in the Reproducibility paragraph of Sec.~IV-D.
\changed
Sec.~IV-D.

% =====================================================================
\section{Reviewer 2}

\itemhead{R2-1}{Baselines limited to SVM variants}
\resp
Random forest and XGBoost added throughout; see AE-4 and R1-5. The reviewer's
framing, that the real competitor in industry is rarely an SVM, is correct
and is now reflected in the result: at the main operating point XGBoost, not the
quantum kernel, is the top-ranked model.
\changed
All results sections.

\itemhead{R2-2}{Barren plateaus have a kernel analogue: exponential
concentration of the kernel matrix}
\resp
Named, cited, and, we hope usefully, measured. Exponential concentration is
now stated as background fact (F3) with the reference the reviewer intends, and
it is also the mechanism behind the paper's sharpest new boundary. Fitting
$s(n)\propto n^{-\alpha}$ to the off-diagonal dispersion over $n=4,\dots,10$
gives $\alpha_{\mathrm{ZZ}}=0.590$ against $\alpha_{\mathrm{Z}}=0.292$ at $K=20$,
a ratio of $2.02$, matching the $\binom{n}{2}$ versus $n$ growth in the number
of phase terms; at $K=80$, where the collapse occurs, the ratio is $3.32$, so
the counting argument predicts the ordering of the two rates everywhere and
their ratio only where the kernel does not collapse, and the manuscript now
says so at every place the ratio is quoted. Across the same widths the
dispersion correlates with macro-$F_1$ at $r=+0.77$ for the ZZ kernel (95\% CI
$[+0.04,+0.96]$ on seven widths) and $r=+0.32$ for the Z kernel
($[-0.57,+0.87]$); the contrast between the two is not resolvable at seven
points ($z=0.97$, $p=0.33$), and the manuscript no longer presents it as
evidence that the relation is specific to the concentrating map.
\changed
Sec.~III (F3), Sec.~V-G, Fig.~8, Sec.~VII-B.

\itemhead{R2-3}{$N_{\text{train}}=1000$ with five seeds is a thin statistical
base}
\resp
Accepted. All experiments now use ten runs with fixed seeds. Every comparison is
paired within run and reported with a $t$-based $95\%$ interval, a Wilcoxon
signed-rank test, the paired effect size $d_z$, and the fraction of runs with a
positive difference; $p$-values are Holm-corrected within baseline families.
Training-set size is now a swept variable rather than a fixed choice, from
$N=100$ to $N=10^{4}$.

We also removed a statistical defect the reviewer did not see: the submitted
version certified the entanglement effect with a two-sample $z$-test that
treated two arms sharing a training subset as independent, discarding the
pairing. The revision compares within run.
\changed
Sec.~IV-B, Sec.~IV-C, Sec.~III-D.

\itemhead{R2-4}{Regime map oversells: positive regimes get full statistics,
negative regimes only qualitative treatment}
\begin{revcomment}
``three positive regimes are given full statistical treatment while the two
negative regimes get comparatively brief qualitative explanation $\ldots$
without a matching effect-size/CI treatment.''
\end{revcomment}
\resp
Corrected, and the correction changed the balance of the paper. All $110$ cells
now carry the same five quantities ($\Delta$, $95\%$ CI, Holm-corrected $p$,
$d_z$, and the positive-run fraction) whether the verdict is favourable to
the quantum kernel or against it. The result is $21$ favourable, $21$ against,
$68$ inconclusive; the negative regimes, given equal treatment, turn out to be
as numerous as the positive ones.

The clearest of them: under feature perturbation, six of six comparisons favour
the classical model with effect sizes beyond $2.8$. The manuscript presents this
as the paper's clearest negative result rather than in a closing caveat.
\changed
Sec.~V-C, Sec.~VI, Fig.~9; \texttt{regime\_map\_rows.csv}, all $110$ rows
independently re-derived by \texttt{audit\_c4.py}.

\itemhead{R2-5}{Reference [15]: article number 116990F should be 116990B}
\resp
Verified independently before correcting: the SPIE Digital Library URL contains
\texttt{/116990B/}, the NASA ADS bibcode is \texttt{2021SPIE11699E..0BP}, and
DOI 10.1117/12.2593297 matches. Corrected.
\changed
Bibliography, entry [15].

\itemhead{R2-6}{Reference [26] could not be located; looks fabricated or badly
misattributed}
\resp
We searched twice and independently, by author and title and by volume and
page range, and could not locate the item either. We have removed it, and we
record the removal in the released bibliography file together with the original
string, so that the change is auditable rather than silent.

Rather than simply delete the sentence it supported, we replaced the citation
with a survey that does exist and that we verified before using: A.~Kaissar,
A.~B. Nassif and A.~Bouridane, ``Enhancing network intrusion detection with
quantum machine learning: A comprehensive survey of methods, metrics, and
applications,'' \emph{Future Internet}, vol.~18, no.~5, art.~234, 2026, doi:
10.3390/fi18050234. It fills the role the withdrawn reference was serving and,
being a 2026 survey, also addresses part of AE-2. Two of its findings are
relevant enough to state in our Sec.~II-B: most published QML--IDS approaches
rely on simulated quantum environments and legacy datasets, and evaluation
practice is inconsistent across studies.
\changed
Bibliography; Sec.~II-B.

% =====================================================================
\section{Reviewer 3}

Reviewer~3 recommends rejection on grounds of limited technical novelty. We take
the criticism seriously and do not dispute its factual basis: this paper
introduces no new quantum kernel, feature map, training method, hardware
implementation, or theorem establishing a new regime of advantage. We respond by
saying what we think the contribution is instead, and by reporting what changed.

\itemhead{R3-1}{The model is a standard ZZFeatureMap configuration; technical
novelty is limited}
\begin{revcomment}
``it does not introduce a new quantum kernel, a new feature map, a new training
method, a hardware implementation, or a theoretical result establishing a new
regime of advantage.''
\end{revcomment}
\resp
Correct as stated, and we have not attempted to manufacture novelty of apparatus.
We offer three things instead, and ask that they be weighed on their own terms.

\emph{An ordering that reverses with sample size.} This is not a re-measurement
of a known quantity. To our knowledge no study of quantum kernels on tabular
data sweeps the training-set size, and doing so shows that the answer to
``does the quantum kernel win'' is not a fact about the kernel but a function of
$N$: second-worst of seven models at $N=100$, best at $N=10^{4}$, reversing in a
narrow interval, in both tuning arms. A benchmark conducted at a single $N$,
as every study we are aware of is, including our own submitted version, can
report either outcome and be internally consistent.

\emph{A selection rule that transfers.} The three-stage rule returns
$n^{\ast}=4$ on NSL-KDD and, with no threshold changed, $n^{\ast}=6$ on
UNSW-NB15, recovered on $10/10$ independent subsets. The submitted version
presented $n=4$ as a result; it is properly a rule whose output depends on the
data.

We would add that the question itself is being treated as a legitimate research
question elsewhere. The study Reviewer~4 pointed us to (Carducci, ICAD 2026)
opens by arguing that malware analysts ``do not need another general claim
that quantum computing is either broadly superior or broadly impractical'' but
rather ``need to know which kinds of malware representations justify quantum
effort'', and answers it by locating a bounded intermediate window. We reached the same shape of answer independently, on a
different task and with a different primitive. That two studies in adjacent
domains converge on ``bounded window'' rather than on ``advantage'' is, we
think, a result about the technology rather than a shortfall of either paper.

\emph{A measured mechanism with a quantitative form.} This is the closest we can
honestly come to the reviewer's ``theoretical result establishing a regime'',
and we present it as an empirical regularity rather than a theorem. The
off-diagonal dispersion of the Gram matrix decays as $n^{-\alpha}$ with
$\alpha_{\mathrm{ZZ}}/\alpha_{\mathrm{Z}}=2.02$ at $K=20$, the ratio predicted
by counting $\binom{n}{2}$ pairwise phase terms against $n$ single-qubit ones,
and $3.32$ at $K=80$, where the counting argument predicts the ordering but not
the rate; the dispersion tracks end-task macro-$F_1$ across widths at $r=+0.77$
(95\% CI $[+0.04,+0.96]$). It tells us
where the method stops working, and that boundary is the part of this paper most
likely to hold beyond these two datasets.
\changed
Sec.~I, Sec.~II-D, Sec.~V-D, Sec.~V-F, Sec.~V-G.

\itemhead{R3-2}{The advantage rests on a marginal improvement ($0.854$ vs
$0.838$)}
\resp
The claim has been withdrawn, and the specific comparison the reviewer names no
longer holds in our own numbers: with a tree ensemble included, XGBoost attains
$0.8503$ against QSVM-ZZ's $0.8469$, with overlapping intervals. We report this
as an ordering of point estimates, not a separation. The reviewer's underlying
point, that a difference of this size does not support an advantage claim,
is one the revision now makes itself, quantitatively: with ten runs and Holm
correction, the smallest difference reaching any verdict in the map is $0.010$
macro-$F_1$, and differences as large as $0.031$ remain inconclusive.
\changed
Sec.~V-B, Sec.~VI.

\itemhead{R3-3}{The propositions are well-established facts and their framing as
proved propositions is inappropriate}
\resp
Agreed, and corrected. Propositions~1, 3 and~4 are now background facts
(F1)--(F3) with citations and no proofs. Proposition~2, which is specific to the
comparison this paper makes, is retained as Lemma~1, and, as
Sec.~\ref{sec:selfreported} reports, we found in the course of doing so that its
stated coefficients were wrong. It is restated correctly with a full proof in
Appendix~A.
\changed
Sec.~III-B, Appendix~A, erratum subsection.

\itemhead{R3-4}{All experiments use very few qubits and small data, yet are
framed as NISQ-aware}
\resp
Three changes, and one thing we did not do.

We extended the width: the study now includes a full re-run at $K=80$, where the
selection rule requires $n^{\ast}=8$ qubits and $112$ two-qubit gates, and the
UNSW-NB15 analysis operates at six qubits. The result is negative and we report
it as the paper's sharpest boundary: at eight qubits all $48$ paired comparisons
favour the classical model, and the entanglement ablation reverses sign. Widening
the circuit does not extend the advantage; it removes it, for a reason we can
measure in the Gram matrix.

We extended the data: training-set size now runs to $N=10^{4}$ with results
reported on the complete $22{,}544$-row test partition.

We added a backend-derived noise model with gate, readout and relaxation errors.

What we did not do is execute on a quantum processor. Rather than leave
``NISQ-aware'' ambiguous, the manuscript now defines it: a statement about
circuit budget (qubit count, depth and two-qubit gate count chosen under an
explicit cost model) and not a claim of hardware validation. Sec.~VII-A states
which effects are consequently outside our evidence. The execution path for a
real backend is in the released code and has been verified in dry run; we report
no hardware numbers and make no claim about them.
\changed
Sec.~V-G, Fig.~8, Sec.~VII-A, Assumption~1.

\itemhead{R3-5}{Results are already established in prior benchmarking studies
(arXiv:2403.07059, arXiv:2409.04406)}
\resp
We have read both and cite both, and our aggregate finding agrees with theirs:
general-purpose feature maps rarely beat well-tuned classical baselines. Table~I
positions this work against them explicitly.

We would note, respectfully, that neither of those studies introduces a new
kernel, feature map, training method, hardware implementation, or theoretical
result either; both are benchmarks. We raise this not to argue that they should
not have been published (we think they should) but because the criterion in
R3-1, applied uniformly, would exclude them, which suggests the criterion for a
benchmark is not novelty of apparatus but whether it changes what the field
believes.

On that criterion the relevant question is what our study does that theirs does
not. We read all four in full before answering, and Table~I is built from those
full texts rather than from abstracts. Bowles \emph{et al.} vary model family
across $160$ mostly synthetic datasets at up to $18$ qubits; Schnabel and Roth
vary kernel type across $64$ datasets at up to $15$ qubits; Cirillo \emph{et al.}
vary qubit count from $2$ to $10$ and plot performance against gate count;
Gillani \emph{et al.} run a $4/6/8/12$-qubit ablation. Circuit width is
therefore \emph{not} an axis we can claim as ours, and we no longer do.

\emph{None of the four varies the training-set size.} Each fixes it at a single
value: fixed per-task subsamples in the first, $M=240$ in the second, a $15\%$
class-balanced subsample in the third, the full training partition in the
fourth. That is the variable that determines the answer here, and it is the one
axis on which this work is not duplicated. We do not claim breadth otherwise;
Gillani \emph{et al.} evaluate four datasets to our two, at eight qubits to our
four and six, under a stricter error-rate criterion.
\changed
Sec.~II-C, Sec.~II-D, Table~I.

% =====================================================================
\section{Reviewer 4}

Reviewer~4 found a false theorem and a selection step that filters nothing. Both
are confirmed. We are grateful, and we have tried to meet the standard the
review sets by auditing the rest of the paper to the same depth and reporting
what that turned up, including in the one theoretical statement we retained.

\itemhead{R4-1}{Related work: Carducci, ICAD 2026, on the same question for
malware detection}
\resp
Cited and discussed, in a subsection we added for it. We are grateful for the
pointer: the parallel is closer than we expected, and it changed how we frame
our own boundary.

Carducci reports that the measured advantage is confined to an
\emph{intermediate} representational window (classical methods remain
stronger both below and above it) and that the pattern survives parameter
tuning, tested specifically to rule out a favourable-setting artefact. All three
features recur in our results on a different task: our favourable regime is
bounded from below by training-set size and from above by circuit width; we
report both a frozen and a re-tuned arm for exactly the reason that paper
re-tunes; and the representations where that work finds quantum discovery useful
are those dominated by local pairwise structure, which is what the
$\binom{n}{2}$ pairwise phase terms of the ZZ map encode.

We must be clear about the evidentiary status. We were unable to obtain the full
text through our institutional access. The comparison rests on the abstract, the
opening of the introduction, and the index terms, from which we can establish
that the evaluation is a retrieval setting scored by Recall@$k$ and that kernel
methods are among the techniques involved, but not the qubit count nor whether
the training-set size is varied. We say exactly this in Sec.~II-D, present the
parallel as suggestive rather than as evidence, and do not tabulate the study in
Table~I, whose cells are otherwise read from full texts. If the reviewer has
access and finds that our reading misstates the work, we would be glad to be
corrected.
\changed
New Sec.~II-D; bibliography; Table~I caption states the exclusion and why.

\itemhead{R4-2}{Theorem~1 appears misstated; and the Pareto step may not filter
anything}
\begin{revcomment}
``It says $\tilde F(4) > \tilde F(3)$ but the opposite is shown in Table III
$\ldots$ that may mean Pareto might not be doing any filtering at all, so it
might not be a contribution after all.''
\end{revcomment}
\resp
Both observations are correct.

On the theorem: the proof asserted $\tilde F(4)>\tilde F(3)$ while Table~III of
the same submission reports $\tilde F(4)=0.471 < \tilde F(3)=0.628$. With the
correct values the conclusion does not follow: at
$\alpha=\beta=\gamma=1/3$ the objective is maximised at $n=2$ $(J=0.551)$ and
decreases monotonically to $0.059$ at $n=10$, exactly as the reviewer observed
from the table.

On Pareto: $V(n)$ increases monotonically while $\tilde F(n)$ decreases and
$Q(n)$ increases, so no candidate dominates any other and every candidate is
Pareto-optimal. The step selected nothing.

We removed Theorem~1 and the objective $J$ rather than repairing them, because
the second point shows the construction was not doing the work attributed to it.
The replacement is lexicographic and has no weights: retain $n$ with
$V(n)\ge0.85$; among those keep $n$ with $\mathrm{KTA}(n)\ge0.95\,\mathrm{KTA}_{\max}$;
among those take the smallest $Q(n)$. On NSL-KDD this returns $n^{\ast}=4$; on
UNSW-NB15, unchanged, $n^{\ast}=6$.
\changed
Sec.~III-C (new Definition~1 replacing Definition~4, Algorithm~1, Eq.~(6) and
Theorem~1), erratum subsection, Fig.~1 replacing the Pareto frontier figure.

\itemhead{R4-3}{Reproducibility is asserted but no machinery is provided}
\resp
A public repository is now cited by URL and commit hash, containing the
implementation, protocol files, all intermediate artifacts, the figure scripts,
and four audit scripts. We regard the audits as the substantive part of this
response: \texttt{audit\_c4.py} recomputes every published statistic from raw
per-run records using an independent implementation rather than calling the
pipeline's own statistics functions, since using the code that produced a number
to check that number lets shared errors through.
\changed
Sec.~IV-D.

\itemhead{R4-4}{Proposition~3 is cited to another paper but not derived}
\resp
Correct. The former Proposition~3 restated a known bound and its ``proof'' added
nothing. It is now background fact (F2), cited to the source and used only
diagnostically, with no proof.
\changed
Sec.~III-B.

\itemhead{R4-5}{No numbers are given for the rare-attack subset, so the
$+6.7$-point claim cannot be verified}
\begin{revcomment}
``I don't know if I see any number for the rare-attack subset, so it is not
possible to verify the paper's claim that `At N=500 QSVM-ZZ still leads by +6.7
points over SVM-RBF on the rare-attack subset, with a Cohen's d of +0.68'\,''
\end{revcomment}
\resp
The claim cannot be verified because it cannot be reproduced. On inspecting the
code that produced the submitted results we found that no rare-class metric was
computed at any point. We withdraw the claim.

The revision reports rare-class metrics directly, at every training-set size and
for every model. At $N=10^{4}$ under the natural prior the rare-class $F_1$ is:

\begin{center}
\small
\begin{tabular}{@{}lcccccc@{}}
\toprule
SVM-RBF & QSVM-ZZ & QSVM-Z & XGBoost & Random forest & Poly-2 & Linear \\
\midrule
0.585 & 0.507 & 0.421 & 0.334 & 0.331 & 0.109 & 0.089 \\
\bottomrule
\end{tabular}
\end{center}

Two things follow, and we report both. The quantum kernel detects rare attacks
far better than either tree ensemble, which the macro average conceals and which
matters operationally. But SVM-RBF is better still, by $0.078$, so the honest
statement, and the one the manuscript makes, is that an RBF-kernel SVM is the
strongest rare-attack detector in this study. Signed per-sample decision margins
are also reported, as the submitted version's Cohen's $d$ claim implied but did
not deliver.
\changed
Sec.~V-E; \texttt{c4\_rare\_attack\_natural.csv}.

\itemhead{R4-6}{The quantum kernel's hyper-parameters were not tuned}
\resp
They are now, by the identical procedure applied to every classical model; see
R1-3. Because the quantum kernel does not depend on $C$, its grid costs only
repeated fits on an already-computed Gram matrix, so there was never a
computational reason for the asymmetry.
\changed
Sec.~IV-B.

% =====================================================================
\section{Corrections we identified ourselves}
\label{sec:selfreported}

Auditing the manuscript to the standard Reviewer~4's comments set turned up
three further defects that no reviewer raised. We report them because a reviewer
comparing the manuscript against the released artifacts would find them, and
because two of them are of the same kind as the errors that were raised.

\textbf{(a) Table~III mislabelled its third column.} The caption defined
$\tilde F(n)=1/\mathrm{DBI}(n)$, but the tabulated values are the ANOVA $F$
statistic; the inverse Davies--Bouldin index takes the values $1.143$, $0.982$
and $0.922$ at $n=2,3,4$. The same column carried a third, also incorrect,
label in an earlier draft. No conclusion in the revision depends on the
quantity, since it has been removed along with $J$, but the mislabelling was
carried across versions and we record it.

\textbf{(b) The former Proposition~2 had incorrect coefficients.} This is the
one theoretical statement the revision retains, and it was wrong. It asserted
\[
  K^{\mathrm{Q}} = 1 - r^{2}\!\sum_i\Delta_i^2
  - r^{2}\!\!\sum_{i<j}\!\Delta\varphi_{ij}^2 + O(\|\cdot\|^4)
  \quad\text{for all } r .
\]
Two things are wrong. The coefficients are not equal: the correct weights are
$1$ and $\tfrac14$, which follows from the gate decomposition, since
$P(2x_i)=\exp(-\mathrm{i}x_iZ_i)$ up to global phase while the
CNOT--$R_Z(\varphi_{ij})$--CNOT trio realises
$\exp(-\mathrm{i}\tfrac{\varphi_{ij}}{2}Z_iZ_j)$. And the $r^{2}$ prefactor does
not hold at all: for $r\ge2$ the interleaved Hadamard layers make the relevant
operator non-diagonal, and no coefficient pair makes the two-invariant form
correct.

Numerically, against an exact statevector kernel verified against Qiskit to
$1.3\times10^{-15}$: at $r=1$ a least-squares fit recovers exactly $1$ and
$\tfrac14$ with $R^2=1.000000$ at $n=3,4,6$, and the corrected expansion's
residual decays as $\|\cdot\|^{3.9}$, reaching $5\times10^{-11}$ at
$\|\cdot\|=10^{-3}$; with the submitted coefficients the residual decays only as
$\|\cdot\|^{2.0}$, so the claimed second-order accuracy is not attained. At
$r=2$, the operating point of this study, the two-invariant form attains
only $R^2=0.60$ to $0.94$ with coefficients that change sign between expansion
points, while the unrestricted quadratic form still attains $R^2=0.999996$.

The lemma is restated for $r=1$ with the correct coefficients and proved in full
in Appendix~A, with an explicit remark on why it does not extend. The
substantive consequence drawn from it, that SVM-Poly2 is the sharpest
classical comparator, is unaffected, since it depends on which monomials
appear and not on their coefficients. The verification is released as
\texttt{verify\_lemma1.py}.

\textbf{(c) The training subsets used in the submitted version were
rare-enriched.} The per-run training frames contain $10\%$ rare-attack records
against a natural rate of $0.83\%$, an enrichment of roughly twelvefold that
was not stated in the submitted manuscript. This is the direct cause of the
withdrawn low-data claim: under enrichment there is no crossover, and under the
natural prior there is. Both sampling regimes are now reported side by side and
the enrichment is stated in the protocol.

% =====================================================================
\section{Note on the marked-up manuscript}

A marked copy accompanies this submission. New or changed text is set on a
yellow background: the abstract, the body of every subsection that is new,
rewritten or corrected with respect to the submitted version, and every
bibliography entry that was added or corrected. Because essentially every
paragraph of the manuscript is new text, most of the copy is yellow; to make
the marking informative, each subsection heading additionally carries one of
four tags (\textsc{new}, \textsc{rewritten}, \textsc{corrected},
\textsc{retained}) with one sentence stating what changed and the item in this
letter that prompted it. Subsections tagged \textsc{retained} are left plain.

The marked copy also lists, on its cover page, the material \emph{removed} from
the submitted version, which by its nature cannot be shown in the body:
Theorem~1, Definition~4, Algorithm~1, the objective $J(n)$, three tables and
four figures. Both files are generated from a single source and differ in one
flag, so the marked copy and the clean copy cannot drift apart.

% =====================================================================
\section{Changes to the bibliography (for the Editor-in-Chief and the Associate Editor)}

Per the TETC policy on references, every change to the bibliography with
respect to the submitted version is listed here with its reason. The submitted
version had $36$ entries; the revision has $38$. Added and corrected entries
are set on a yellow background in the reference list of the marked copy. No
self-citation was added or removed; the manuscript contains none.

\emph{Six entries added.} All six are cited in the body and discussed, not
merely listed.
\begin{itemize}[leftmargin=1.3em,itemsep=2pt,topsep=3pt]
  \item Cirillo, Esposito and Seo, \emph{Quantum Machine Intelligence} 8(1):27,
    2026. Named by Reviewer~1 (R1-1) as a missing benchmark; now the closest
    prior study in Table~I and Sec.~II-B.
  \item Bowles, Ahmed and Schuld, arXiv:2403.07059, 2024, and Schnabel and
    Roth, arXiv:2409.04406, 2024. Both named by Reviewer~3 (R3-5) as evidence
    that our results were already established; discussed in Sec.~II-C.
  \item Carducci, \emph{Proc. IEEE Int. Conf. AI and Data Analytics (ICAD)},
    2026. Named by Reviewer~4 (R4-1); discussed in Sec.~II-D.
  \item Gillani \emph{et al.}, arXiv:2608.18155, 2026, a calibration- and
    noise-aware benchmark of QML for intrusion detection on the same datasets,
    found by the authors while answering R1-1 and R3-5; without it Table~I
    would be incomplete.
  \item Kaissar, Nassif and Bouridane, \emph{Future Internet} 18(5):234, 2026,
    a survey of quantum methods for intrusion detection. It takes the role the
    unverifiable former [26] held, and its DOI (10.3390/fi18050234) was checked
    to resolve to the item.
\end{itemize}

\emph{Four entries removed.}
\begin{itemize}[leftmargin=1.3em,itemsep=2pt,topsep=3pt]
  \item Former [26], Rahman, Kayed, Aljahdali and Ali, ``Quantum machine
    learning for cybersecurity: a comprehensive review,'' attributed to
    \emph{IEEE Access} 2023. Reviewer~2 (R2-6) could not locate it; neither
    could we, by author and title or by volume and page range. Removed, and the
    removal is recorded in the released bibliography file.
  \item Former [34]--[36], Platt 1999, Naeini \emph{et al.} 2015, Guo
    \emph{et al.} 2017. These supported the probability-calibration material
    of the submitted version, which was withdrawn together with the advantage
    claims it accompanied; the revision does not discuss calibration, and no
    sentence cites them.
\end{itemize}

\emph{One entry corrected.} [15], Payares and Mart\'inez-Santos, SPIE 11699:
article number 116990F corrected to 116990B (DOI 10.1117/12.2593297), as
Reviewer~2 pointed out (R2-5).

\emph{Affiliations, acknowledgments, funding.} No co-author has changed
affiliation since the original submission. The funding statement (University of
Danang - University of Science and Technology, Project T2026-02-10TN) is
unchanged. Author biographies are included in the revised manuscript, as the
decision letter requests.

% =====================================================================
\section{Summary}

\begin{center}
\small
\begin{tabular}{@{}lcccccc@{}}
\toprule
& AE & R1 & R2 & R3 & R4 & Total \\
\midrule
Items raised                  & 6 & 10 & 6 & 5 & 6 & 33 \\
Addressed with new results    & 6 &  9 & 6 & 4 & 5 & 30 \\
Addressed in part, stated as such & 0 & 1 & 0 & 1 & 1 & 3 \\
\bottomrule
\end{tabular}
\end{center}

\vspace{4pt}
The three items we describe as addressed only in part, and what remains
outstanding in each:

\begin{itemize}[leftmargin=1.3em,itemsep=2pt,topsep=3pt]
  \item \textbf{R4-1}: Carducci (2026) is cited and positioned, but we were
    unable to obtain the full text through our institutional access, so our
    discussion of it rests on the abstract and the bibliographic record.
  \item \textbf{R1-5}: we added a random forest and XGBoost, and deliberately
    did not add CatBoost, TabNet or FT-Transformer. We took the alternative the
    comment offers and softened the practical claims accordingly.
  \item \textbf{R3-4}: we extended the study to eight qubits, added a
    backend-derived noise model, and defined what ``NISQ-aware'' claims. We did
    not execute on a quantum processor, and we make no hardware claim.
\end{itemize}

We thank the reviewers again. The paper reports a narrower and partly negative
result than the one we submitted, and it is a better paper for it.

\vspace{10pt}
Sincerely,\\[4pt]
The authors\\
\emph{Corresponding author:} Quan Tran Anh Vo

\end{document}
